考虑恒等式

\begin{equation}\tag{4.5.31}
Z(t,V/\mathbb F_q)
=\exp\!\left(
\sum_{\nu\geq1}
\operatorname{card}V(\mathbb F_{q^\nu})\frac{t^\nu}{\nu}
\right)
\end{equation}

及

\begin{equation}\tag{4.5.32}
\det\!\left(1-t\alpha_B^{\operatorname{ord}(q)}\right)
=\exp\!\left(
-\sum_{\nu\geq1}
\operatorname{tr}\!\left((\alpha_B)^{\nu\operatorname{ord}(q)}\right)
\frac{t^\nu}{\nu}
\right),
\end{equation}

公式 (4.5.27) 可改写为

\begin{equation}\tag{4.5.33}
Z(t,V/\mathbb F_q)\prod_{i=0}^{n+1}(1-q^it)
=\prod_{j=1}^{n+2}
\left(
\prod_{\substack{B\subset S\\\operatorname{card}(B)=j}}
\Delta_B(t)
\right)^{-(-\delta)^j\varphi^{n+2-j}}.
\end{equation}

\medskip
\noindent\emph{底本注.}
印本将 (4.5.23) 编号为 (4.5.2.3), 在 (4.5.25) 的第一个指数中
将 $\operatorname{card}(A)$ 缩写为 $\operatorname{car}(A)$,
并在最后的 zeta 恒等式中再次使用编号 (4.5.23).
校订法文层和英文层使用上列连续编号及完整的基数记号;
照录法文层保留印本形式.

\subsection*{4.6.}
现研究级数 $\Delta_B$ 的整性. 置

\begin{equation}\tag{4.6.1}
\mathcal O=\text{$K$ 的整数环},
\end{equation}

并用下式定义 $1+t\mathcal O[[t]]$ 上的滤过

\begin{equation}\tag{4.6.2}
F^i(1+t\mathcal O[[t]])
=1+q^it\mathcal O[[q^it]]
=\varphi^i(1+t\mathcal O[[t]]).
\end{equation}

有

\begin{align}
\varphi(F^i)&\subset F^{i+1}, \tag{4.6.3}\\
\delta(F^i)&\subset F^i, \tag{4.6.4}\\
\delta&\text{ 在下列空间上诱导恒等映射: }F^i/F^{i+1}. \tag{4.6.5}
\end{align}

\subsubsection*{引理 4.6.6.}
有 $\Delta_B(t)\in1+t\mathcal O[[t]]$; 此外, 若整数 $j$ 满足

\begin{equation}\tag{4.6.7}
jd<\operatorname{card}(B),
\end{equation}

则

\begin{equation}\tag{4.6.8}
\Delta_B(t)\in F^j(1+t\mathcal O[[t]]).
\end{equation}

